\section{結言}
今回のDaily Xでは、Astraそのものの新規発表より、モデルをどの規模で訓練し、どれだけ効率よく利用し、現実の専門作業へどう適用するかが主役となった。Pachockiのエッセイは、その急速な能力拡大と同時に、alignmentやmonitoringの未解決性を正面から論点化している。

また、Vals AIのコスト性能主張や用途別比較、Free Pause TokensやCROCODILの研究共有は、単純なモデルサイズ競争だけでなく、推論効率、評価設計、モデル間協調へ競争軸が広がっていることを示した。半導体供給の噂も含め、モデル、計算資源、利用制約、研究方法論が同じ観測窓の中で密接に結び付いた一日だった。
